\scanpage{7}
\uline{\unsure{Démontrons le|Démonstration du} Lemme} 1 \quad Il sera une conséquence du

\uline{Lemme 2} \inserted{(Mumford)} Soit $\deleted{\unsure{A}}B$ schém. abélien sur un corps \unsure{[alg. clos ?]|(alg. clos ?)|(alg. clos)} $k$, et $L$ un Module inv. sur $B$ alg. équivalent à zéro, \deleted{\unsure{Al|Alors}} et non isomorphe à $\mathcal O_B$. Alors on a
\[
\deleted{\unsure{-}}\quad H^i(X,L)=0\quad\deleted{\unsure{\text{si }i>0}}\quad\text{pour \unsure{\deleted{tout}|tout} }i
\]
\[
\deleted{\text{et }\ H^0(X,L)=
\begin{cases}
0&\text{si }L\unsure{\not\simeq|\nsim}\underline O_B\\
k&\text{si }L\simeq\underline O_B
\end{cases}}
\]

\deleted{\uline{Démonstration du lemme 2} L'assertion sur le $H^0$ est triviale, et il reste à prouver que $H^i(X,L)=0$ si $i\unsure{>}$.}

Nous allons ici admettre le lemme 2, et allons montrer comment en \deleted{conclure} \inserted{déduire} le lemme 1. Nous voyons déja grâce à lemme 2 que $Rp_*(W)$ est concentré sur la section unité de $B^*$. Nous allons définir un homomorphisme canonique (où on pose $e_{B^*}=e$)
\[
Rp_*(W)\longrightarrow\deleted{\unsure{e_*(\underline O_S)[-n]}}e_*(\underline{\mathcal O}_S)\otimes_{\underline O_S}\tau_{B^*}[-n]
\]
ou, ce qui revient au $=$, \inserted{(comme $n$ est la dim relative de $B/S$, donc $R^ip_*(W)=0$ si $i>n$)} un homomorphisme canonique
\[
R^np_*(W)\longrightarrow\deleted{\unsure{e_*(\underline O_S)|\otimes_S}}e_*(\tau_{B^*})\deleted{\unsure{[-n]}}
\]
ou encore un hom.
\[
\deleted{e_{B}}\quad e^*(\mathbb{R}^n p_*(W))\longrightarrow\deleted{\unsure{\underline O_S|\otimes_S}}\tau_{B^*}
\]
\diagram{\begin{tikzcd}[ampersand replacement=\&]
B^*\times_S B\arrow[d,"p"] \& B\arrow[l,"{e'}"']\arrow[d,-,"q"]\\
B^* \& S\arrow[l,"e"']
\end{tikzcd}}
Or en dim. $n$ on a \inserted{formule de} changement de base, et \inserted{comme $e'^*(W)\simeq\underline O_B$,} on est ramené à définir un hom.
\[
R^n q_*(\underline O_B)\longrightarrow\deleted{\unsure{\underline O_S|\otimes_S}}\tau_{B^*}
\]
Or on sait que l'on a un isom canonique
\[
R^nq_*(\underline{\mathcal O}_B)\simeq\textstyle\bigwedge^{n}R^1q_*(\underline O_B)\simeq\bigwedge^{n}t_{B^*}
\]
où $t_{B^*}$ est le faisceau tangent de $B^*$ le long de $e$. On \deleted{\unsure{trouve donc}} \inserted{prend alors} \uline{l'hom identique} $\bigwedge^n t_{B^*}\simeq\tau_{B^*}$.
% archive sheet number: 6
